\begin{lemma}[Wild odd upper residual coefficients]
\label{mod:cases-wildodd-upperresidualcoefficients}
Using the High/Low normalization certificate supplied by
UpperResidualNormalization and its common/source coordinate data, prove the
natural upstairs polar coefficient, upper/source and upper/lower coordinate
transport, inverse-Frobenius display, and every upstairs polar row.  Retain the
genuine quadratic contribution and use HigherSymmetricVanishing to prove that
degrees $3,\ldots,p-1$ and the terminal cancellation create no hidden affine
contribution.  Then prove every upstairs affine row and record all permitted
representative translations.  At even conductor parity, prove directly from
the source-tied PhaseReduction functions that the displayed High and Low rows
are constant one and therefore have coefficient pair $(0,0)$.

For the actual High row, also expose the exact source bridge.  An element of
exact signed order has nonzero initial residue.  If
$a=(s+1)-2d$ (with no assumption that $a\geq 0$), transport the norm initial
residue through the residue-degree-one equivalence and prove
\[
 \overline{N(u)/\varpi_F^a}
   =\left(e^{-1}\!\left(\overline{u/\varpi_K^a}\right)\right)^p
\]
without a uniformizer residue scalar.  Construct the canonical displayed
source unit from the literal displacement $u x$ by
\texttt{positiveUnitOfLattice} at depth
$((s+1)+2)/2=(s+3)/2$, and record that its field value is exactly $1+ux$;
the residue variable in $x$ is the prescribed inverse-Frobenius source
coordinate.

Finally identify the actual norm-generator polar coefficient with $\eta_0$
and the zero-index source/base polar coefficient with
$\eta=\eta_0d_0$.  The latter proof must retain the stationary numerator as
its quotient-class representative until the exact ratio is rewritten as
$A\,n$, use the signed initial-residue identity above and the integral
\texttt{wildOddCompactEta} formula, and introduce no additional scalar.  These
identifications imply that the named boundary ratio is the same $d_0$ used by
the actual upper row.

This node does not assemble the final PhasePreparation package, prove the
lower coefficient table, assemble the final public coefficient table, or
perform finite-field phase cancellation.
\LeanFile{LanglandsFirstMainLemma/Cases/WildOdd/UpperResidualCoefficients.lean}
\lean{LanglandsFirstMainLemma.wildOdd_upperResidualCoefficients}
\leanok
\SourceText{Lemmas lem:critical-polar-coordinate, lem:critical-polynomial-coordinate, lem:odd-upper-higher-vanishing, and lem:residual-parameter-transport; Proposition prop:odd-coefficient-table, upper rows.}
\uses{mod:cases-wildodd-upperresidualnormalization,mod:cases-wildodd-highervanishing}
\FormalizationContract
This is one cohesive upper-table responsibility, expected to occupy roughly
1800--4600 Lean lines.  Upper residual normalization, higher-symmetric
estimates, and public assembly must remain in their existing separate nodes.
Do not split the polar and affine rows merely to impose a smaller artificial
line limit; line count alone does not determine the architecture.  In the
High source bridge, do not replace the literal $ux$ unit by a compact residue
surrogate, reuse a cancelled old-layer parameter after a conductor drop, or
discard a representative translation from the Lamprecht phase.
\end{lemma}
